\section{The Corner of the World: The Mission to the English (596--601)}

\subsection{The Gallic slave market and the first design (595--596)}

Among all the enterprises of Gregory's pontificate, none reshaped Western
Christendom more permanently than the mission to the Anglo-Saxon kingdoms of
Britain. When the story is stripped of later legend and read first from
Gregory's own correspondence, it is even more striking---because it shows a
pope watching northern Europe with hardheaded pastoral realism while Lombards
raided outside Rome's gates.

The earliest surviving reference to the English in Gregory's register appears
in September 595, in a brief business letter to Candidus, a Roman presbyter
dispatched to administer the Church's small patrimony in Gaul (\work{Ep.}\
VI.7). Because Gallic gold coins could not be spent at full value in Italy,
Gregory told Candidus to spend the revenues locally on two works of mercy:
clothing for the poor, or ``English boys of about seventeen or eighteen years
of age, who may profit by being given to God in monasteries.'' Since the boys
available for purchase in the Gallic markets were pagans, Gregory added a
concrete pastoral precaution: ``I desire that a presbyter be sent hither with
them to provide against the case of any sickness occurring on the way, that he
may baptize those whom he sees to be about to die'' (\work{Ep.}\ VI.7). Behind
that brief instruction lies a stark sixth-century reality: Gregory knew the
English through the western slave trade and had already conceived the plan of
redeeming English youths to form them as monks and future missionaries.

Within months he moved to a direct mission. In the spring of 596 he chose
Augustine, provost (\latin{praepositus}) of his own monastery of Saint Andrew
on the Caelian Hill, to lead a band of monks north toward Kent, where King
Ethelbert had married the Christian Frankish princess Bertha. Yet the
enterprise nearly collapsed before reaching the English Channel. Reaching
southern Provence---where Augustine visited the Bishop of Aix and the famous
island monastery of L\'erins---the monks grew terrified by reports of the
fierce, barbarous people whose language they did not know, and sent Augustine
back to Rome to ask Gregory's permission to turn home (Barmby, note to
\work{Ep.}\ VI.51, citing Bede, \work{Historia Ecclesiastica} I.23, and John
the Deacon, \work{Vita} II.33).

Gregory refused to let fear abort the mission. On 23 July 596 he sent Augustine
back to his waiting companions with an invigorating letter appointing Augustine
their abbot:

\begin{quote}
Since it had been better not to have begun what is good than to return back
from it when begun, you must, most beloved sons, fulfil the good work which
with the help of the Lord you have begun. Let, then, neither the toil of the
journey nor the tongues of evil-speaking men deter you; but with all instancy
and all fervour go on with what under God's guidance you have commenced,
knowing that great toil is followed by the glory of an eternal reward. \dots\
May Almighty God protect you with His grace, and grant to me to see the fruit of
your labour in the eternal country; that so, even though I cannot labour with
you, I may be found together with you in the joy of the reward; for in truth I
desire to labour. (\work{Ep.}\ VI.51)
\end{quote}

At the same time, Gregory supplied the practical diplomatic scaffolding the
monks had lacked on their first departure. He armed Augustine with commendatory
letters to the bishops of Marseilles, Tours (or Turni), Arles, Vienne, Autun,
and Aix, to Abbot Stephen of L\'erins, to Arigius the Patrician of Gaul, to
the young Frankish kings Theodoric of Burgundy and Theodebert of Austrasia,
and to their grandmother Queen Brunehild (\work{Ep.}\ VI.52--59). Writing to
the Frankish rulers, Gregory gave his own plain statement of why Rome had
intervened: ``it has come to our knowledge that the nation of the Angli is
desirous, through the mercy of God, of being converted to the Christian faith,
but that the priests in their neighbourhood neglect them, and are remiss in
kindling their desires by their own exhortations'' (\work{Ep.}\ VI.58; cf.\
VI.59). He instructed Augustine to take Frankish presbyters along as
interpreters, ensuring the Roman monks would not land mute on Kentish shores.

\subsection{First harvest and pastoral strategy (597--601)}

Supported by the Frankish courts---including King Clotaire II of Neustria
(Barmby, Prolegomena, citing \work{Ep.}\ XI.61)---Augustine and his monks
landed on the Isle of Thanet in Kent early in 597 (Bede, \work{Historia
Ecclesiastica} I.25, V.24, as reported by Barmby). The astonishing speed of
their reception stands verified in Gregory's own hand barely a year later. In
July 598, writing from a sickbed to Patriarch Eulogius of Alexandria, Gregory
shared his exultation over a harvest at the far rim of the known world:

\begin{quote}
For while the nation of the Angli, placed in a corner of the world, remained up
to this time misbelieving in the worship of stocks and stones, I determined,
through the aid of your prayers for me, to send to it, God granting it, a monk
of my monastery for the purpose of preaching. And he, having with my leave been
made bishop by the bishops of Germany, proceeded, with their aid also, to the
end of the world to the aforesaid nation; and already letters have reached us
telling us of his safety and his work; to the effect that he and those that
have been sent with him are resplendent with such great miracles in the said
nation that they seem to imitate the powers of the apostles in the signs which
they display. Moreover, at the solemnity of the Lord's Nativity which occurred
in this first indiction, more than ten thousand Angli are reported to have been
baptized by the same our brother and fellow-bishop. (\work{Ep.}\ VIII.30)
\end{quote}

Gregory's reference to Augustine's consecration ``by the bishops of Germany''
(\latin{a Germanis episcopis}) reflects loose Roman geography for Gaul; Bede
(\work{Historia Ecclesiastica} I.27) and John the Deacon (\work{Vita} II.36,
reported by Barmby) specify Archbishop Virgilius of Arles. The baptism of more
than ten thousand English at Christmas 597---even allowing for the rounded
rhetoric of an enthusiastic mission report---confirmed that Kent had opened its
doors to the Gospel.

A second wave of letters belongs to June and July 601 (Book XI of the
\work{Registrum}, summarized in Barmby's Prolegomena from Migne PL 77 and Bede,
\work{Historia Ecclesiastica} I.27--32).
When Augustine sent the priest Laurence and the monk Peter back to Rome to seek
reinforcements and answers to practical pastoral dilemmas, Gregory dispatched
Mellitus, Justus, Paulinus, and Rufinianus carrying the metropolitan pallium
for Augustine, liturgical books, altar vessels, and relics. With them went a
remarkable dossier: a letter to King Ethelbert (\work{Ep.}\ XI.66) and another
to Queen Bertha (\work{Ep.}\ XI.29); organizational blueprints dividing Britain
into two future metropolitan provinces at London and York, each with twelve
suffragans (\work{Ep.}\ XI.65); the famous \work{Libellus Responsionum}
(\work{Ep.}\ XI.64) answering Augustine's questions on liturgy, marriage,
charity, and relations with the older British bishops; and a warning to
Augustine himself not to grow proud over reported miracles (\work{Ep.}\ XI.28).

One final letter of July 601 reveals Gregory's mature missionary wisdom at its
finest. After instructing Ethelbert to destroy pagan shrines, Gregory reflected
further and sent an urgent message after Abbot Mellitus to overtake the party
in Gaul (\work{Ep.}\ XI.76, reported in Barmby, Prolegomena). Rather than
tearing down well-built pagan temples, the missionaries should destroy only the
idols, sprinkle the buildings with holy water, set up altars with relics, and
convert the traditional seasonal feasts into Christian celebrations honoring
the martyrs---leading a stubborn people step by step rather than by violent
leaps. That same flexible principle---holding unity of faith while refusing
needless uniformity in local custom---had already guided his counsel to Leander
of Seville ten years earlier, assuring him that differing local usages do not
injure the Church so long as the one faith remains intact (\work{Ep.}\ I.43).

Nowhere in all these contemporary letters does Gregory mention an encounter
with fair-haired English boys in the Roman forum or the celebrated puns
\latin{Non Angli sed Angeli}. That beloved story first surfaces more than a
century later in English tradition, where its origin and meaning require
separate assessment below.
